\documentclass{article}

\textwidth=6.5in
\textheight=8.5in
\voffset=-54pt
\oddsidemargin=0in
\evensidemargin=0in
\setlength{\parskip}{8pt}
\setlength{\parindent}{0pt}

\usepackage{amsmath, amssymb}
\usepackage{ifthen}

\usepackage[inline]{enumitem}
\setlist{topsep=0pt, parsep=8pt}

\usepackage[rgb,table]{xcolor}\convertcolorsUfalse
\usepackage{graphicx}

% change hyperref options
\PassOptionsToPackage{hyphens}{url}
\usepackage[bookmarks,colorlinks,linkcolor=black,citecolor=blue,pdfstartview=FitH,urlcolor=blue]{hyperref}
\hypersetup{pdfpagemode=UseNone}
\usepackage{bookmark}

\usepackage{graphicx}
\usepackage{multicol}
\usepackage{framed}
\usepackage{array}

\usepackage{icomma}

\usepackage{pifont}

\usepackage{soul}

\usepackage{tikz}
\usetikzlibrary{
  matrix,%
  % enables the snake paths
  decorations.pathmorphing,%
  % enables bigger arrows
  decorations.markings,%
  decorations.pathreplacing,%
  decorations.text,%
  calc,%
  patterns,%
  shapes.geometric,%
  arrows,
  decorations.shapes,
  positioning,
  plotmarks,
  intersections
}
\tikzset{>=latex} % sets default arrow type in tikz
\tikzset{font=\normalsize}
\tikzset{mark size=1.5pt, mark options=thin}
\tikzset{pin distance=4pt,
  every pin edge/.style={<-, thin, shorten <= -2pt}}

\interfootnotelinepenalty=10000
\renewcommand{\thefootnote}{(\arabic{footnote})}

\usepackage{multirow, bigdelim}

% change to \usepackage[sols]{handout} to see solutions
\usepackage{handout}

\newcommand\iscurrentenumi[1]{%
  \ifthenelse{\equal{#1}{\theenumi}}{}{#1}}
\setenumerate[1]{ref=\#\arabic*}
\setenumerate[2]{ref=\protect\iscurrentenumi{\theenumi}(\alph*)}
\newcommand\enumiiprefix[2]{%
  \ifthenelse{{\equal{#1}{\theenumi}}\AND{\equal{#2}{(\alph{enumii})}}}{}{}%
  \ifthenelse{{\equal{#1}{\theenumi}}\AND{\NOT{\equal{#2}{(\alph{enumii})}}}}{#2}{}%
  \ifthenelse{\equal{#1}{\theenumi}}{}{#1#2}%
}
\setenumerate[3]{ref=\protect\enumiiprefix{\theenumi}{(\alph{enumii})}\roman*}

\definecolor{purple}{rgb}{0.5,0,1.0}

\usepackage{amsthm}

% make URLs print in the same font
\usepackage{url}
\urlstyle{same}

\usepackage{pifont}
\def\exend{\hfill\ding{118}}

%\makeatletter\@addtoreset{equation}{section}\makeatother
%\renewcommand{\theequation}{\arabic{section}.\arabic{equation}}

\newtheoremstyle{theorem}% name
{8pt}%      Space above
{0pt}%      Space below
{\itshape}%         Body font
{}%         Indent amount (empty = no indent, \parindent = para indent)
{\bfseries}% Thm head font
{.}%        Punctuation after thm head
{.5em}%     Space after thm head: " " = normal interword space;
                                %       \newline = linebreak
{}%         Thm head spec (can be left empty, meaning `normal')

\theoremstyle{theorem}

\let\Theorem\undefined
\newtheorem{Theorem}[equation]{Theorem}
\let\Definition\undefined
\newtheorem{Definition}[equation]{Definition}
\let\Summary\undefined
\newtheorem{Summary}[equation]{Summary}
\let\Conjecture\undefined
\newtheorem{Conjecture}[equation]{Conjecture}
\let\Fact\undefined
\newtheorem{Fact}[equation]{Fact}

\renewcommand{\thesubsection}{\arabic{subsection}}
\makeatletter
\def\@seccntformat#1{\@ifundefined{#1@cntformat}%
   {\csname the#1\endcsname\quad}%       default
   {\csname #1@cntformat\endcsname}}%    enable individual control
\newcommand\section@cntformat{}
\makeatother

\usepackage{fancyhdr}
\lhead{\textcolor{white}{This content is protected and may not be shared, uploaded, or distributed.}}
\rhead{}
\rfoot{\vspace*{\fill}\footnotesize\textcopyright{} \emph{Harvard Math Ma}}
\renewcommand{\headrulewidth}{0pt}
\pagestyle{fancy}
\begin{document}

  \htitle{Exponentials \& Logs Practice Session}

\begin{enumerate}

\item \emph{The basics}

  \begin{enumerate}
  \item
    \textbf{The Definition of Logarithms.} $\log_b x$ is\ldots

    \pspace{0.35in}

    Saying $y = \log_b x$ is equivalent to what exponential equation?  

    \pspace{0.35in}

  \item Make up a few examples;

    $\log_{10} \left( \fitb{0.5in}{} \right)= \fitb{0.5in}{}$ \qquad $\displaystyle \log_{10} \left( \fitb{0.5in}{} \right) = \fitb{0.5in}{}$ \qquad $\displaystyle \log_{10} \left( \fitb{0.5in}{} \right) = \fitb{0.5in}{}$

    \pspace{0.35in}

  \item
    \begin{problem}[0.35in]
      What does $\ln A$ mean?
    \end{problem}

  \item
    \begin{problem}[\vfill]
      Write down the exponent rules.
    \end{problem}

  \item
    \begin{problem}[\nstretch{0.75}]
      Write down the log rules.
    \end{problem}

  \item
    \begin{problem}[\vfill]
      On a single set of axes, sketch graphs of $e^x$, $3^x$, $\ln x$, and $\log_3 x$.
    \end{problem}

  \item
    \begin{problem}[\stretch{0.5}]
      $b^x$ and $\log_b x$ are inverse functions. Write down two identities that this fact implies.
    \end{problem}

  \item $\dfrac{d}{dx} (e^{kx}) = $ \hfill $\dfrac{d}{dx} (b^x) = $ \hfill $\dfrac{d}{dx} (\ln x) = $ \hfill $\dfrac{d}{dx} (\log_b x) = $ \hfill \ 

    \pspace{\nstretch{0.5}}

  \end{enumerate}

\item Find the critical points of:

  \begin{enumerate}
  \item
    \begin{problem}[\vfill]
      $f(x) = 3^{2x + 1} \cdot 4^{x - 2} - \dfrac{x}{\ln 2}$
    \end{problem}

  \item
    \begin{problem}[\vfill\newpage]
      $g(x) = \dfrac{\ln(3x^5)}{2} - e^{2 \ln x}$
    \end{problem}

  \end{enumerate}

\item
  Anna and Beatrice are studying two different radioactive substances; both are decaying exponentially. 
  \begin{itemize}[nosep]
  \item Anna's substance currently has a mass of 10 g, and it has a half life of 50 days.
  \item Beatrice's substance currently has a mass of 8 g, and it decays by 10\% every 10 days. 
  \end{itemize}

  \begin{enumerate}
  \item
    \begin{problem}[\vfill]
      Your classmate has forgotten what ``half life'' means. Explain it to them. 
    \end{problem}

  \item
    \begin{problem}[\vfill]
      After 50 days, will Beatrice's substance have decayed by exactly 50\%, more than 50\%, or less than 50\%? How do you know?
    \end{problem}

  \item
    \begin{problem}[\vfill]
      Does Anna's substance or Beatrice's substance have a longer half life? How do you know?
    \end{problem}

  \item
    \begin{problem}[\nstretch{2.5}]
      Find the half-life of Beatrice's substance. 
    \end{problem}

  \item
    \begin{problem}[\vfill]
      Anna's substance currently has a mass of 10 g. How quickly is its mass changing with respect to time (right now)? Please give an instantaneous rate of change, and include units. 
    \end{problem}

  \item
    \begin{problem}[\nstretch{2}]
      Is there a time when Anna's substance and Beatrice's substance have the same mass? If so, find that time; if not, why not? 
    \end{problem}

  \end{enumerate}

\item % based on data from  https://ecos.fws.gov/ecp/report/species-listings-by-year-totals

  Over the past few decades, the number of endangered species in the U.S. has been growing. A reasonable model for the number of endangered species in the U.S. $t$ years after 1990 is $f(t) = 1000 \log_7 (t + 5) - 200$.

  \begin{enumerate}

  \item
    \begin{problem}[\vfill]
      According to the model, was the number of endangered species in the U.S. in 1990 greater than or less than 300? 
    \end{problem}

  \item
    \begin{problem}[\vfill]
      Sketch a graph of $f(t)$. Please label any asymptotes, and label 1 point on the graph with exact integer coordinates. 
    \end{problem}

  \item
    \begin{problem}[\nstretch{0.25}]
      Do you think the model gives a good estimate of the number of endangered species in 1985? 
    \end{problem}

  \item
    \begin{problem}[\vfill]
      According to the model, will the number of endangered species ever be 1800? If so, when? 
    \end{problem}

  \item
    \begin{problem}[\vfill\newpage]
      Imagine that, in another country, the number of endangered species is modeled by $g(t) = 1000 \log_7 (t + 47) - 1200$. Is there a time when the U.S. and this other country have the same number of endangered species? If so, when?
    \end{problem}

  \end{enumerate}

\end{enumerate}

\end{document}
